\documentclass[12pt,a4paper,oneside]{article}
\usepackage[brazil]{babel}
\usepackage[utf8]{inputenc}
\usepackage{olymp}


\setcounter{problem}{4}

\begin{document}

\begin{problem}{Meia volta volver!}{volver}{3}

O sargento chega e os soldados logo se colocam em fila. Mas a fila não fica bem formada! Uns estão virados pra um lado e outros pra outro. Veja um exemplo com 8 soldados, em que uns estão virados pra direita (\texttt{D}) e outros pra esquerda (\texttt{E}).

\vspace{-6pt}
\begin{center}
    \texttt{DEEEDDDE}
\end{center}

Insatisfeito com a falta de coordenação dos soldados o sargento ordena que, ao seu sinal, todos que estão frente a frente deem meia volta. No exemplo acima os soldados na $1^a$ e $2^a$ posição e os da $7^a$ e $8^a$ posição (o da esquerda está virado pra direita e o da direita pra esquerda). O problema é que isso gera outros conflitos... veja o que acontece quando o sargento dá o sinal:

\begin{center}
    \texttt{\underline{\textbf{DE}}EEDD\underline{\textbf{DE}}}

    \texttt{EDEEDDED}
\end{center}

Agora os soldados nas posições 2 e 3 estão de frente um pro outro. O mesmo acontece com os das posições 6 e 7. O sargento dá outro sinal e esses soldados dão meia-volta, chegando a:

\vspace{-6pt}
\begin{center}
    \texttt{EEDEDEDD}
\end{center}

Mais um sinal é dado e chegam a:

\vspace{-6pt}
\begin{center}
    \texttt{EEEDEDDD}
\end{center}

Ainda mais um:

\vspace{-6pt}
\begin{center}
    \texttt{EEEEDDDD}
\end{center}

Embora não estejam todos virados para o mesmo lado o sargento se dá por satisfeito, afinal não há nenhum soldado frente a frente (olhando um pro outro).

Sua tarefa neste problema é escrever o formato da fila a cada sinal do sargento, até que não haja soldados frente a frente.


%\Notes
%
%Observações, se tiver.

\InputFile

A entrada contém duas linhas. A primeira contém um número inteiro $N$, o número de soldados. A segunda contém $N$ caracteres `\texttt{D}' ou `\texttt{E}', indicando a direção inicial de cada soldado. Restrições: $2 \leq N \leq 100$.


\OutputFile

Seu programa deve gerar uma ou mais linhas na saída. A primeira é a formação inicial dos soldados, ou seja, é igual à formação fornecida na entrada. A última é a formação final dos soldados, quando nenhum mais está de frente pro outro. As demais linhas indicam as formações intermediárias, a cada sinal do sargento.

Note que todos as linhas possuem $N$ caracteres. Note ainda que pode haver caso em que a saída contém somente uma linha. Isso acontecerá caso a formação inicial não tenha soldado de frente pra outro, ou seja, já é uma formação final.


% \Notes
% Preste atenção aos tipos das variáveis, pois $F(90) \approx 2^{61}$.

\Example

\begin{example}
\exmp{
8
DEEEDDDE
}{
DEEEDDDE
EDEEDDED
EEDEDEDD
EEEDEDDD
EEEEDDDD
}%
\end{example}

\begin{example}
\exmp{
7
DEEEEED
}{
DEEEEED
EDEEEED
EEDEEED
EEEDEED
EEEEDED
EEEEEDD
}%
\end{example}

\begin{example}
\exmp{
5
EEEDD
}{
EEEDD
}%
\end{example}

\begin{example}
\exmp{
4
EEEE
}{
EEEE
}%
\end{example}


\begin{example}
\exmp{
10
DEDEDEDEDE
}{
DEDEDEDEDE
EDEDEDEDED
EEDEDEDEDD
EEEDEDEDDD
EEEEDEDDDD
EEEEEDDDDD
}%
\end{example}

%Comentários sobre o exemplo, se necessário.

\end{problem}

\end{document}
